\chapter{Methodology}
\label{ch:method}

This chapter describes the approach taken to meet the objectives in \cref{sec:objectives}. It first gives the overall workflow, then formulates the prediction problem, describes the teacher (RouteNet-Fermi) and the proposed student (Queue-Link Fusion), and finally sets out how the student is trained with knowledge distillation.

\section{Overall Workflow}
\label{sec:workflow}

\Cref{fig:workflow} summarises the workflow. The same data pipeline turns each simulated network scenario into a graph, and four models are trained from it. The \emph{baseline} is RouteNet-Fermi trained on the dataset's original training split, which contains only 10-node topologies. The \emph{teacher} has the same architecture but is trained on a re-split of the data that includes topologies of 10 to 170 nodes. The \emph{student} is trained twice on the same re-split: once on the ground truth alone (\emph{student-scratch}) and once with distillation from the frozen teacher (\emph{student-KD}). All four models are then evaluated on one held-out test set of 200- to 300-node topologies, and their size and inference latency are measured.

\begin{figure}[ht]
\centering
\begin{tikzpicture}[
  node distance=6mm and 8mm,
  box/.style={draw, rounded corners=2pt, align=center, font=\small, minimum height=10mm, text width=27mm, fill=gray!8},
  data/.style={draw, align=center, font=\small, minimum height=10mm, text width=27mm, fill=blue!6},
  arr/.style={-{Stealth[length=2.2mm]}, thick}
]
\node[data] (raw) {All-mixed dataset\\(simulated scenarios)};
\node[box, right=of raw] (hg) {Graph construction\\(paths, links, queues)};
\node[data, right=of hg] (split) {Topology-based\\data splits};

\node[box, below=12mm of hg] (teach) {Teacher\\RouteNet-Fermi\\(10--170-node train)};
\node[box, left=8mm of teach] (base) {Baseline\\RouteNet-Fermi\\(10-node train)};
\node[box, right=8mm of teach] (stud) {Student\\Queue-Link Fusion\\(scratch / KD)};

\node[data, below=12mm of teach] (test) {Held-out test set\\(200--300 nodes)};
\node[box, below=6mm of test, text width=70mm] (eval) {Evaluation: MAPE, parameters, model size, latency};

\draw[arr] (raw) -- (hg);
\draw[arr] (hg) -- (split);
\coordinate (bus) at ($(split.south)+(0,-6mm)$);
\draw[thick] (split.south) -- (bus);
\draw[thick] (base.north |- bus) -- (stud.north |- bus);
\draw[arr] (base.north |- bus) -- (base.north);
\draw[arr] (teach.north |- bus) -- (teach.north);
\draw[arr] (stud.north |- bus) -- (stud.north);
\draw[arr, dashed] (teach) -- (stud);
\draw[arr] (base) |- (test);
\draw[arr] (teach) -- (test);
\draw[arr] (stud) |- (test);
\draw[arr] (test) -- (eval);
\end{tikzpicture}
\caption{Overall workflow of the project. The dashed arrow shows the teacher's predictions being used as an additional training signal for the student.}
\label{fig:workflow}
\end{figure}

\section{Problem Formulation}
\label{sec:problem}

A network scenario consists of a topology, a routing configuration and a traffic matrix. The topology has a set of directed links $\mathcal{L}$; each link $l$ has a capacity $c_l$ and inherits a scheduling policy $\pi_l \in \{\text{WFQ}, \text{SP}, \text{DRR}, \text{FIFO}\}$ from its source node. Each link has a set of output queues $\mathcal{Q}(l)$, with at most three queues per link. Queue $q$ has a buffer size $b_q$, a priority index $r_q$ and, for weighted policies, a normalised weight $w_q$.

Traffic is described by a set of flows $\mathcal{P}$, each following a fixed path. A path $p$ is an ordered sequence of links $(l_{p,1}, \dots, l_{p,|p|})$ and, at each hop $k$, uses the queue $q_{p,k}$ that its type-of-service (ToS) class maps to on that link. Each flow has a mean bit rate $\lambda_p$, a packet generation rate $n_p$, a traffic model $m_p$ (one of seven, encoded one-hot) and up to eight parameters of its time distribution $\bm{\theta}_p$.

The task is to predict, for every flow $p$ in a scenario, its mean end-to-end delay $y_p$ as measured by the packet-level simulator that generated the dataset. A model $f$ receives the whole scenario and outputs $\hat{y}_p$ for all flows at once. Accuracy is measured with the mean absolute percentage error,
\begin{equation}
\mathrm{MAPE} = \frac{100}{|\mathcal{P}|} \sum_{p \in \mathcal{P}} \frac{|y_p - \hat{y}_p|}{|y_p|},
\label{eq:mape}
\end{equation}
which is also the training loss for the baseline, the teacher and the scratch student. A relative error is used because delays in the dataset span more than two orders of magnitude (\cref{sec:datasets}), so an absolute error would be dominated by the slowest flows.

Two derived quantities are computed from the inputs before any learning takes place. The offered load of a link is
\begin{equation}
\rho_l = \frac{1}{c_l} \sum_{p \,:\, l \in p} \lambda_p ,
\label{eq:load}
\end{equation}
and the mean packet size of a flow is $s_p = \lambda_p / n_p$.

\section{Graph Representation}
\label{sec:graph}

Each scenario is converted into a graph whose nodes are of three types: \emph{paths}, \emph{links} and \emph{queues}. A path is connected to every link and every queue it traverses; a queue is connected to the link it belongs to. This representation follows RouteNet-Fermi \cite{ferriolgalmes2023routenetfermi} and is built by the data pipeline distributed with its reference implementation, which reads the simulator output through the dataset's API and produces, for each scenario, the feature tensors and index tensors needed for message passing (for example, which links and queues each path visits, in order).

The graphs are large. One scenario from the dataset's original test split, used during development to inspect the pipeline, contained 852 links, 2,124 queues and 57,360 flows, and therefore 57,360 delay labels. Every training step in this project processes one such scenario; there is no batching across scenarios.

\section{Teacher Model: RouteNet-Fermi}
\label{sec:teacher}

The teacher is RouteNet-Fermi \cite{ferriolgalmes2023routenetfermi}, used without architectural changes. It keeps a hidden state $\mathbf{h}_p$, $\mathbf{h}_l$ and $\mathbf{h}_q$ of dimension 32 for every path, link and queue.

\paragraph{Initial states.}
Each state is initialised by a two-layer MLP with ReLU activations applied to the entity's features, where $z(\cdot)$ denotes z-score normalisation with fixed dataset statistics:
\begin{align}
\mathbf{h}_p^{(0)} &= \mathrm{MLP}_p\big([\,z(\lambda_p),\, z(n_p),\, \mathrm{onehot}(m_p),\, z(\bm{\theta}_p)\,]\big) && \text{(17 inputs)}\\
\mathbf{h}_l^{(0)} &= \mathrm{MLP}_l\big([\,\rho_l,\, \mathrm{onehot}(\pi_l)\,]\big) && \text{(5 inputs)}\\
\mathbf{h}_q^{(0)} &= \mathrm{MLP}_q\big([\,z(b_q),\, \mathrm{onehot}(r_q),\, w_q\,]\big) && \text{(5 inputs)}
\end{align}

\paragraph{Message passing.}
The model then runs $T = 8$ iterations, each consisting of three stages.
\begin{enumerate}
  \item \emph{Links and queues to paths.} A GRU \cite{cho2014gru} reads, hop by hop, the concatenated states of the queue and link used at each hop, starting from the current path state:
  \begin{equation}
  \mathbf{s}_p^{(k)} = \mathrm{GRU}_p\big(\,[\mathbf{h}_{q_{p,k}} \,\|\, \mathbf{h}_{l_{p,k}}],\ \mathbf{s}_p^{(k-1)}\big), \qquad \mathbf{s}_p^{(0)} = \mathbf{h}_p,
  \end{equation}
  and the new path state is $\mathbf{h}_p \leftarrow \mathbf{s}_p^{(|p|)}$. The intermediate states $\mathbf{s}_p^{(k)}$ are kept.
  \item \emph{Paths to queues.} Each queue sums the states of the paths that use it, taken at the point where each path arrives at the queue, and updates its state with a second GRU cell:
  \begin{equation}
  \mathbf{h}_q \leftarrow \mathrm{GRU}_q\Big(\sum_{(p,k)\,:\,q_{p,k}=q} \mathbf{s}_p^{(k-1)},\ \mathbf{h}_q\Big).
  \end{equation}
  \item \emph{Queues to links.} Each link runs a third GRU over the states of its queues, starting from its own state, and keeps the final output as its new state.
\end{enumerate}

\paragraph{Readout.}
After the last iteration, an MLP (32--16--16--1) maps each per-hop path state to a value $o_{p,k}$ interpreted as the queue occupancy seen by flow $p$ at hop $k$. Delay is then computed from the physics of a store-and-forward network rather than predicted directly:
\begin{equation}
\hat{y}_p = \underbrace{\sum_{k=1}^{|p|} \frac{o_{p,k}}{c_{l_{p,k}}}}_{\text{queueing delay}} \;+\; \underbrace{s_p \sum_{k=1}^{|p|} \frac{1}{c_{l_{p,k}}}}_{\text{transmission delay}} .
\label{eq:readout}
\end{equation}
Dividing by link capacity means that the network only has to learn occupancy, which is less sensitive to the absolute link speed than delay itself. This is one of the reasons RouteNet-Fermi is reported to generalise to networks with link capacities outside the training range \cite{ferriolgalmes2023routenetfermi}.

In total the teacher has 27,025 trainable parameters, and each forward pass performs $8 \times 3 = 24$ message-passing stages.

\section{Proposed Student Model: Queue-Link Fusion}
\label{sec:student}

The student is designed to cut the cost of message passing while keeping the parts of RouteNet-Fermi that are most closely tied to how delay arises: the per-hop recurrent reading of each path and the capacity-based readout of \cref{eq:readout}.

\subsection{Design idea}

In a real router, the queues of a link are not independent objects; they are buffers on that link's output port, and their combined behaviour determines how long a packet waits before transmission on that link. The student exploits this by merging each link and its queues into a single \emph{interface} entity before message passing begins. The tripartite path--link--queue graph of the teacher becomes a bipartite path--interface graph (\cref{fig:graphs}). This removes the queue states, the queue GRU, the paths-to-queues stage and the queues-to-links stage.

\begin{figure}[ht]
\centering
\begin{subfigure}[b]{0.48\textwidth}
\centering
\begin{tikzpicture}[
  ent/.style={circle, draw, minimum size=8mm, font=\small, inner sep=0pt},
  every edge/.style={draw, -{Stealth[length=1.6mm]}, thin}
]
\node[ent, fill=orange!20] (p1) at (0,2.2) {$p_1$};
\node[ent, fill=orange!20] (p2) at (2.4,2.2) {$p_2$};
\node[ent, fill=green!15] (q1) at (0,0.9) {$q_1$};
\node[ent, fill=green!15] (q2) at (1.2,0.9) {$q_2$};
\node[ent, fill=green!15] (q3) at (2.4,0.9) {$q_3$};
\node[ent, fill=blue!12] (l1) at (0.6,-0.4) {$l_1$};
\node[ent, fill=blue!12] (l2) at (2.4,-0.4) {$l_2$};
\draw[-{Stealth[length=1.6mm]}] (q1) -- (p1);
\draw[-{Stealth[length=1.6mm]}] (q2) -- (p2);
\draw[-{Stealth[length=1.6mm]}] (q3) -- (p2);
\draw[-{Stealth[length=1.6mm]}] (q1) -- (l1);
\draw[-{Stealth[length=1.6mm]}] (q2) -- (l1);
\draw[-{Stealth[length=1.6mm]}] (q3) -- (l2);
\draw[-{Stealth[length=1.6mm]}, bend right=40] (l1) to (p1);
\draw[-{Stealth[length=1.6mm]}, bend left=10] (l1) to (p2);
\draw[-{Stealth[length=1.6mm]}, bend right=40] (l2) to (p2);
\end{tikzpicture}
\caption{Teacher: paths, links and queues (3 stages per iteration).}
\end{subfigure}
\hfill
\begin{subfigure}[b]{0.48\textwidth}
\centering
\begin{tikzpicture}[
  ent/.style={circle, draw, minimum size=8mm, font=\small, inner sep=0pt}
]
\node[ent, fill=orange!20] (p1) at (0,2.2) {$p_1$};
\node[ent, fill=orange!20] (p2) at (2.4,2.2) {$p_2$};
\node[ent, fill=blue!12, minimum size=11mm] (i1) at (0.6,0.0) {$i_1$};
\node[ent, fill=blue!12, minimum size=11mm] (i2) at (2.4,0.0) {$i_2$};
\node[font=\scriptsize, align=center] at (0.6,-1.0) {$l_1 + \{q_1,q_2\}$};
\node[font=\scriptsize, align=center] at (2.4,-1.0) {$l_2 + \{q_3\}$};
\draw[{Stealth[length=1.6mm]}-{Stealth[length=1.6mm]}] (i1) -- (p1);
\draw[{Stealth[length=1.6mm]}-{Stealth[length=1.6mm]}] (i1) -- (p2);
\draw[{Stealth[length=1.6mm]}-{Stealth[length=1.6mm]}] (i2) -- (p2);
\end{tikzpicture}
\caption{Student: paths and fused interfaces (2 stages per iteration).}
\end{subfigure}
\caption{Graph structures used by the teacher and the student, illustrated for two flows crossing two links. In the student, each interface $i$ represents a link together with all of its queues.}
\label{fig:graphs}
\end{figure}

\subsection{Fusing queues into links}

For each link $l$, the features of its queues are averaged and appended to the link's own features. Writing $\overline{(\cdot)}_l$ for the mean over $q \in \mathcal{Q}(l)$, the initial interface state is
\begin{equation}
\mathbf{h}_l^{(0)} = \mathrm{MLP}_I\Big(\big[\,\rho_l,\ \mathrm{onehot}(\pi_l),\ \overline{z(b_q)}_l,\ \overline{\mathrm{onehot}(r_q)}_l,\ \overline{w_q}_l,\ |\mathcal{Q}(l)|/3\,\big]\Big),
\label{eq:fusion}
\end{equation}
an 11-dimensional input. The averaged one-hot priority vector records which priority levels exist on the link, and the last term records how many queues it has. The fusion is computed once per forward pass with segment-mean operations and adds no trainable parameters.

\subsection{Message passing and readout}

The student uses $T = 4$ iterations of two stages and a state dimension of 16:
\begin{enumerate}
  \item \emph{Interfaces to paths:} $\mathbf{s}_p^{(k)} = \mathrm{GRU}_p\big(\mathbf{h}_{l_{p,k}},\ \mathbf{s}_p^{(k-1)}\big)$ with $\mathbf{s}_p^{(0)} = \mathbf{h}_p$, and $\mathbf{h}_p \leftarrow \mathbf{s}_p^{(|p|)}$.
  \item \emph{Paths to interfaces:} $\mathbf{h}_l \leftarrow \mathrm{GRU}_I\big(\sum_{(p,k)\,:\,l_{p,k}=l} \mathbf{s}_p^{(k-1)},\ \mathbf{h}_l\big)$.
\end{enumerate}
The path embedding uses the same 17 inputs as the teacher. The readout MLP is reduced to 16--8--8--1, and the delay is assembled exactly as in \cref{eq:readout}. The student therefore performs $4 \times 2 = 8$ message-passing stages per forward pass, one third of the teacher's 24, and has 4,505 trainable parameters (\cref{tab:params}).

\subsection{What the fusion gives up}
\label{sec:fusion-cost}

The fusion is not free of information loss, and this should be stated explicitly. In the teacher, the path GRU receives the state of the specific queue that the flow's ToS class maps to at each hop, so two flows crossing the same link in different priority classes see different queue states. In the student, every flow crossing a link sees the same interface state, and the path features (\cref{sec:teacher}) do not include the ToS class. The student therefore has no direct way to tell a high-priority flow from a low-priority flow on the same link if their traffic features are similar. Under strict-priority or weighted scheduling these flows can experience very different delays. This limitation is a deliberate trade for speed, and its likely effect on accuracy is discussed with the results in \cref{sec:discussion}.

\section{Knowledge Distillation}
\label{sec:kd}

\subsection{Distillation loss}

The student-KD model is trained with online, response-based distillation \cite{hinton2015distilling,gou2021kdsurvey}. At every training step the frozen teacher is run on the same scenario, and the student's loss combines its error against the simulator ground truth with its disagreement with the teacher:
\begin{equation}
\mathcal{L}_{\mathrm{KD}} = \alpha \cdot \underbrace{\mathrm{MAPE}\big(\mathbf{y}, \hat{\mathbf{y}}^{S}\big)}_{\text{hard loss}} \;+\; (1-\alpha) \cdot \underbrace{\frac{1}{|\mathcal{P}|}\sum_{p\in\mathcal{P}} \big(\hat{y}^{T}_p - \hat{y}^{S}_p\big)^2}_{\text{soft loss (MSE)}} ,
\label{eq:kd}
\end{equation}
where $\hat{\mathbf{y}}^{T}$ and $\hat{\mathbf{y}}^{S}$ are the teacher's and student's predictions and $\alpha = 0.5$. Only the student's parameters are updated.

Because delay is a regression target, there is no class distribution to soften with a temperature, as in the original formulation of distillation. The teacher's per-flow predictions are used directly as soft targets. They differ from the ground truth in a structured way, since they reflect what a model of this family can and cannot represent, and the hypothesis being tested is that matching them gives the student a smoother and easier target than the simulator labels alone.

\subsection{Choice of teacher}

The teacher used for distillation is the topology-diverse model (trained on 10- to 170-node topologies), not the baseline. The reason is that the student is evaluated on larger topologies than either model saw in training, and the teacher's predictions are only a useful target if they themselves extrapolate well. The results in \cref{sec:overall} confirm that the topology-diverse teacher is the more accurate of the two on the held-out set.

\section{Topology-Diverse Training Data}
\label{sec:diverse}

The dataset's original training split contains only 10-node topologies, while its original test split spans 50 to 300 nodes. To give the teacher wider coverage while still keeping a genuinely unseen test range, the data were re-split by topology size: topologies of 50 to 170 nodes were moved from the original test split into the training and validation splits, and only topologies of 200 to 300 nodes were kept for testing. Splitting by whole topology, rather than by individual scenario, ensures that no test topology was seen during training or model selection. The exact procedure and the resulting splits are given in \cref{sec:datasets}.

\section{Summary of Design Choices}

\Cref{tab:design} summarises the main design decisions and the reason for each.

\begin{table}[ht]
\centering
\small
\caption{Main design choices and their rationale.}
\label{tab:design}
\begin{tabularx}{\textwidth}{@{}p{42mm} Y@{}}
\toprule
\textbf{Choice} & \textbf{Rationale} \\
\midrule
RouteNet-Fermi as teacher & Published state-of-the-art GNN for per-flow delay with queue and scheduling support; open reference implementation and dataset. \\
Re-split by topology size & Gives the teacher wider coverage while keeping test topologies (200--300 nodes) fully unseen. \\
Merge queues into links & Removes one entity type, one GRU and two of three message-passing stages; queues are physically part of a link. \\
State dimension 16, 4 iterations & Further reduces parameters and computation; chosen as half of the teacher's settings. \\
Keep capacity-based readout & Preserves the physical structure of delay (\cref{eq:readout}), which supports extrapolation to unseen link speeds. \\
MAPE as hard loss & Relative error suits targets spanning orders of magnitude; same loss as the baseline and teacher. \\
MSE on teacher predictions & Regression target with no class distribution to soften; direct matching of predictions is the simplest response-based option. \\
\bottomrule
\end{tabularx}
\end{table}
